% !TEX root = ../../cookbook.tex
\newpage
\recipe{Ciabatta}{\href{https://www.youtube.com/watch?v=jA95tUK2RZU&t=4s}{ricetta di Brian Lagerstrom}, Per 2 ciabatte grandi.}{}{}

\begin{multicols}{2}
	\subsection*{Ingredienti}
	
	\subsubsection*{Per la biga}
	\begin{ingredients}
		\item 175 g di acqua calda
		\item \sfrac{1}{4} di cucchianino di lievito secco
		\item 225 gdi farina (11.7\% di proteina)
	\end{ingredients}
	
	\subsubsection*{Per l'autolisi}
	\begin{ingredients}
		\item 180 gdi acqua calda
		\item 250 g di farina
	\end{ingredients}
	\columnbreak
	\subsection*{ }
	
	\subsubsection*{Per il mix finale}
	\begin{ingredients}
		\item 40 g di acqua calda
		\item 5 g di lievito secco
		\item 10 g di sale
		\item Tutta la biga
	\end{ingredients}
\end{multicols}

% --- Preparation ---
\textbf{Fare la biga, idealmente il giorno prima:} \\
Mettere \textbf{acqua}, \textbf{lievito}, e \textbf{farina} in un contenitore dai bordi alti, mischiare bene, e tappare.
Lasciar fermentare a temperatura ambiente per 6-24 h.

\textbf{Autolisi:} \\
Mettere \textbf{acqua} e \textbf{farina} in una planetaria. 
Mischiare col gancio a velocità bassa finché ben combinata, poi coprire con uno straccio e lasciar riposare 30 minuti.

\textbf{Mix finale:} \\
Aggiungere al composto dell'autolisi tutti gli \textbf{altri ingredienti} e la \textbf{biga}, mischiare a velocità bassa per 3 minuti finché ben combinati, poi a velocità alta per 5 minuti.
La pasta deve aver pulito i bordi della ciotola ed essere lucida ed elastica. 

\textbf{Fermentazione:} \\
Tasferire il panetto in un contenitore grande ben oleato, coprire e lasciar lievitare 30 minuti.

Eseguire una piega di rinforzo\footnote{Pieghe di rinforzo:  vedi ricetta \nameref{pizza_napolitana} (p. \pageref{pizza_napolitana}), nota numero~\ref{pieghe_rinforzo}.} e lasciar lievitare altri 30 minuti.

Eseguire un'altra piega di rinforzo, oppure, ancora meglio, una \textit{laminatura}\footnote{La laminatura è un processo per sviluppare struttura nel pane, spesso usata in pasticceria con strati di burro per creare molteplici stati individuali, come nei croissants. Nel caso del pane, consiste in stendere la pasta in uno strato sottile su una superficie bagnata, poi piegare i lati in dentro (come una lettera) e infine arrotolarla su se stessa per formare nuovamente un panetto. Brian Lagerstrom dimostra come eseguire una laminatura nel video della ricetta, che trovi in cima a questa pagina, sotto al titolo.}, e lasciar riposare 1 h.

\textbf{Formare le ciabatte:} \\
Infarinare la pasta e la superficie di lavoro, rovesciare la pasta sul piano, staccandola bene dal suo contenitore per non sgonfiarla.
Tagliare la pasta a metà e, \textbf{gentilmente}, senza schiacciarle ne lavorarle ulteriormente, dare alle pagnotte una forma rettangolare, di ciabatta.
Tasferirle su carta da forno infarinata, coprirle, e lasciarle riposare 3o minuti.

\textbf{Cottura:} \\
Pre-riscaldare il forno a 260°C. 
Mettere le ciabatte nel forno insieme ad un contenitore di acqua bollente\footnote{Il contenitore di acqua bollente sostituisce il controllo dell'umidità dei forni da panificio professionali. Ad inizio cottura, il vapore manterrà la superficie del pane umida, rallentando la formazione della crosta e permettendo alla pagnotta di gonfiarsi di piu. Una volta formata la crosta, il vapore la renderà piu croccante e favorirà un colore piu profondo.}, e cuocere per $\sim$25 minuti, finché bella dorata/scura.






% --- Description ---
\begin{recipeblock}
	
\end{recipeblock}
